\documentclass[11pt]{article}
\usepackage[margin=1in]{geometry}
\usepackage{amsmath,amssymb,amsthm}
\usepackage{graphicx}
\usepackage{microtype}
\usepackage{xcolor}
\usepackage[colorlinks=true,linkcolor=blue!55!black,citecolor=blue!55!black,urlcolor=blue!55!black]{hyperref}

\newtheorem{theorem}{Theorem}
\newtheorem{corollary}[theorem]{Corollary}
\newcommand{\TT}{\widetilde T}
\newcommand{\DD}{\widetilde D}

\title{A Schur-normalization proof for the hyperbolic Cauchy transform\\
\large Candidate full solution to Section 3.4(i) of arXiv:math/0204018}
\author{Automated Mathematics Run \texttt{fa\_banach\_001}, lane 2}
\date{19 July 2026}

\begin{document}
\maketitle

\begin{center}
\fcolorbox{blue!55!black}{blue!4}{\parbox{0.89\textwidth}{
\textbf{Status.} Candidate full solution to Open Problem (i), subject to
human verification.  The conclusion applies to the irreducible unitary
realizations specified in the question.  Reducible parameters are covered
sectorwise, with one normalization constant per irreducible sector.}}
\end{center}

\section{Source and exact question}

The source is V.~V.~Kisil, \emph{Spaces of Analytical Functions and
Wavelets---Lecture Notes}~\cite{Kisil}.  It appeared as arXiv:math/0204018
(2002).  In Section 3.4, printed page 47 (PDF page 48), item (i) asks:

\begin{quote}
Demonstrate that Cauchy formula (3.3.23) is an isometry between
$F_2(\TT)$ and $H_\sigma(\DD)$ with suitable norms chosen.  This almost
follows (up to some constant factor) from its property to intertwine two
irreducible representations of $SL_2(\mathbb R)$.
\end{quote}

Figure~\ref{fig:problem} reproduces the source statement.

\begin{figure}[ht]
  \centering
  \includegraphics[width=\textwidth]{figures/open_problem_crop.png}
  \caption{Section 3.4 and the complete text of Open Problem (i) in the
  source paper~\cite{Kisil}.}
  \label{fig:problem}
\end{figure}

The Cauchy transform in (3.3.23) is denoted below by
\[
  W_\sigma:F_2(\TT)\longrightarrow H_\sigma(\DD).
\]
The source's Appendix ``Boundedness of the Singular Integral Operator
$W_\sigma$'' treats the four boundary singularities in the principal-value
sense and states the resulting $L^2$ boundedness.  The general reduced-wavelet
construction proves equivariance, and item (i) explicitly identifies the
relevant representations as irreducible.  These facts are the only
source-specific inputs used below.

\section{Definitions and proof intuition}

For a unitary representation $\pi:G\to\mathcal U(\mathcal H)$, its
\emph{commutant} consists of bounded operators $A$ satisfying
$A\pi(g)=\pi(g)A$ for every $g\in G$.  A bounded map
$W:\mathcal H_1\to\mathcal H_2$ \emph{intertwines} representations $\pi$
and $\rho$ if
\[
  W\pi(g)=\rho(g)W\qquad(g\in G).
\]

\subsection*{Proof Intuition}

Equivariance makes the positive operator $W^*W$ invisible to the group:
it belongs to the commutant of the irreducible source representation.
Every spectral projection of $W^*W$ is then group invariant.  Irreducibility
leaves no nontrivial spectral projection, forcing $W^*W$ to be one positive
scalar times the identity.  Thus $W$ already preserves every norm ratio by
the same constant.  Removing that single constant makes it an isometry.
The same identity makes the range closed; equivariance and target
irreducibility then force surjectivity.

\section{Normalization theorem}

\begin{theorem}[bounded irreducible intertwiners are scalar unitaries]
\label{thm:intertwiner}
Let $G$ be a group, let
$\pi:G\to\mathcal U(\mathcal H_1)$ and
$\rho:G\to\mathcal U(\mathcal H_2)$ be irreducible complex unitary
representations, and let $0\ne W\in\mathcal B(\mathcal H_1,\mathcal H_2)$
satisfy $W\pi(g)=\rho(g)W$ for every $g\in G$.  Then there are a unique
constant $c>0$ and a unitary intertwiner $U:\mathcal H_1\to\mathcal H_2$
such that
\[
  W=\sqrt c\,U.
\]
Equivalently, $W^*W=cI_{\mathcal H_1}$ and
$\|Wf\|_{\mathcal H_2}=\sqrt c\,\|f\|_{\mathcal H_1}$ for all $f$.
\end{theorem}

\begin{proof}
Taking adjoints of the intertwining identity, and using unitarity, gives
\[
  W^*W\pi(g)=W^*\rho(g)W=\pi(g)W^*W
  \qquad(g\in G).
\]
Thus the positive operator $A=W^*W$ commutes with $\pi(G)$.  By bounded
functional calculus, every spectral projection of $A$ also commutes with
$\pi(G)$.  The range and kernel of such a projection are closed invariant
subspaces.  Irreducibility therefore makes every spectral projection either
$0$ or $I$.  Hence the spectrum of $A$ is a singleton and $A=cI$ for some
$c\ge0$.  Since $W\ne0$, we have $c>0$.

It follows that $U=c^{-1/2}W$ is an isometry, so its range is closed.  The
range of $W$ (and hence of $U$) is nonzero and $\rho(G)$-invariant.  Target
irreducibility forces this closed range to equal $\mathcal H_2$.  Therefore
$U$ is a surjective isometry, hence unitary.  Uniqueness of $c$ follows from
$W^*W=cI$.
\end{proof}

\begin{corollary}[solution of Section 3.4(i)]
\label{cor:application}
Fix a parameter $\sigma$ for which the boundary and Hardy realizations in
Section 3.3 of~\cite{Kisil} are the irreducible unitary representations
specified in Open Problem (i), and let $W_\sigma$ be the bounded nonzero
principal-value Cauchy transform (3.3.23).  Then there is a constant
$c_\sigma>0$ such that
\[
  \|W_\sigma f\|_{H_\sigma(\DD)}^2
    =c_\sigma\|f\|_{F_2(\TT)}^2
  \qquad(f\in F_2(\TT)).
\]
Moreover, $W_\sigma$ maps onto $H_\sigma(\DD)$.  With the target norm
\[
  \|h\|_{H_\sigma,\mathrm{iso}}
    :=c_\sigma^{-1/2}\|h\|_{H_\sigma(\DD)},
\]
the original Cauchy formula $W_\sigma$ is an onto isometry.
\end{corollary}

\begin{proof}
The source establishes that $W_\sigma$ is a bounded intertwiner, while the
question specifies irreducibility of the two representations.  The Cauchy
kernel is not identically zero, so the transform is nonzero.  Apply
Theorem~\ref{thm:intertwiner} with $G=SL_2(\mathbb R)$.
\end{proof}

The normalizing constant can be recovered from any single nonzero test vector:
\[
  c_\sigma=\frac{\|W_\sigma f\|^2}{\|f\|^2}.
\]
The theorem proves that this quotient is independent of the chosen $f$.

\section{Even-Clifford sectors}

The source writes parts of the construction using even Clifford-valued
functions.  If that notation displays the Hilbert realization as an orthogonal
sum of irreducible sectors,
\[
  \mathcal H_1=\bigoplus_{j=1}^m\mathcal H_{1,j},\qquad
  \mathcal H_2=\bigoplus_{j=1}^m\mathcal H_{2,j},
\]
and $W_\sigma=\bigoplus_j W_{\sigma,j}$ respects the sectors, the same proof
gives $W_{\sigma,j}^*W_{\sigma,j}=c_{\sigma,j}I$ on every nonzero irreducible
sector.  The invariant target norm
\[
  \|(h_j)_j\|_{\mathrm{iso}}^2
   =\sum_j c_{\sigma,j}^{-1}\|h_j\|^2
\]
makes the full transform isometric.  Thus the Clifford bookkeeping changes
only the number of scalar normalizations, not the argument.

\section{Verification notes and limitations}

\noindent\textbf{1. Proof audit.}  Boundedness makes $W_\sigma^*$ available;
unitarity is used exactly to move group operators through the adjoint;
irreducibility is used twice, first on spectral projections of $W^*W$ and
then on the closed range.  No density-to-surjectivity shortcut is used.

\par\medskip\noindent\textbf{2. Singular-vacuum issue.}  The source notes that its vacuum
vector is not in $L^2$.  The proof above does not use that vector as an
$L^2$ element.  It starts from the principal-value bounded operator asserted
in the source appendix.

\par\medskip\noindent\textbf{3. Reducibility parameters.}  If a value of $\sigma$ makes the
representation reducible, the single-scalar statement need not hold on the
unreduced space.  The sectorwise statement above is the correct conclusion.
Open Problem (i) itself invokes two irreducible representations, and
Corollary~\ref{cor:application} matches precisely that scope.

\par\medskip\noindent\textbf{4. What is not computed.}  The proof determines the normalization
intrinsically but does not evaluate $c_\sigma$ in closed form relative to a
preselected geometric normalization of boundary measure.  The question asks
for suitable norms, so this does not leave a logical gap in the stated
isometry problem.

\par\medskip\noindent\textbf{5. Remaining source questions.}  Section 3.4(ii), concerning
domain deformation via Stokes' theorem, and Section 3.4(iii), concerning a
principal-series Dirac condition, are not addressed.

\section{Bounded novelty check and review recommendation}

The run's solution, attempt, and registry indexes were searched for arXiv id
0204018 and the title/keywords.  Targeted web/arXiv searches through
19 July 2026 used the exact title and combinations of ``Cauchy formula'',
``isometry'', ``principal series'', ``$H_\sigma$'', and ``Kisil''.  They found
the source paper, the author's earlier paper on principal function theory,
and later surveys, but no paper explicitly resolving item (i).  Since the
argument is the standard commutant proof behind Schur's lemma, the claim of
novelty is deliberately modest: this packet records an explicit completion
of the source's stated ``up to a constant'' route.

\textbf{Human-review recommendation: likely valid.}  The decisive source
check is that the selected $\sigma$ gives the positive-definite irreducible
unitary boundary and Hardy realizations claimed in the question, and that the
principal-value transform is nonzero on the selected sector.  Once those
source hypotheses are confirmed, the normalization proof is complete.

\begin{thebibliography}{9}
\bibitem{Kisil}
V. V. Kisil,
\emph{Spaces of Analytical Functions and Wavelets---Lecture Notes},
arXiv:math/0204018 (2002), especially Sections 2.2, 3.3, 3.4 and the appendix
``Boundedness of the Singular Integral Operator $W_\sigma$''.
\end{thebibliography}

\end{document}
